% !PRINTED 625
\weberpage{625}

Since here $b=0$, by (7) and (8) we have $y=2u$, and
from (11) and (15) the values of the $y$ result:
\[
\begin{aligned}
y_1&=y_2=y_3=y_6=-2\sqrt{\alpha}\\
y_4&=(4-2i)\sqrt{\alpha},\quad y_5=(4+2i)\sqrt{\alpha}.
\end{aligned}
\]
From this one obtains for $\beta=0$ the following equation,
identical in $y$:
\weq{17}{\begin{aligned}
&y^6+a_2y^4+a_4y^2+a_6-a_5\sqrt{\Delta}\,y\\
&=(y+2\sqrt{\alpha})^4(y^2-8y\sqrt{\alpha}+20\alpha)\\
&=y^6-20\alpha y^4+240\alpha^2y^2+512\sqrt{\alpha}^5y+320\alpha^3.
\end{aligned}}
But from this one finds the general resolvent for the
equation (13), if according to (16) one puts $\sqrt{\Delta}$ for $-16\sqrt{\alpha}^5$, thus
\weq{18}{y^6-20\alpha y^4+240\alpha^2y^2-32\sqrt{\Delta}\,y+320\alpha^3=0.}

The discriminant $\Delta$ here has, as follows from formula (3)
in \S.\,74, the expression
\weq{19}{\Delta=5^5\beta^4+2^8\alpha^5.}

Simpler still is the equation for $u$, namely
\weq{20}{u^6-5\alpha u^4+15\alpha^2u^2-\sqrt{\Delta}\,u+5\alpha^3=0,}
and if one puts $u^2=v$, there results for $v$ the rational
equation of degree 6
\weq{21}{(v^3-5\alpha v^2+15\alpha^2v+5\alpha^3)^2=\Delta v.}

Another form of this equation is obtained from (17),
if one multiplies the second of the three expressions by the one
obtained from it by changing $\sqrt{\alpha}$ into $-\sqrt{\alpha}$.
In this way one obtains
\[
(v-\alpha)^4(v^2+6\alpha v+25\alpha^2)=0,
\]
and (21) must pass over into this when $\beta=0$ is set.
But since in (21) the part depending on $\beta$ is determined
by (19), the sought form of the resolvent follows:
\weq{22}{(v-\alpha)^4(v^2+6\alpha v+25\alpha^2)=5^5\beta^4v.}

One may further raise the question whether the function $v$
introduced here is always really metacyclic, whether it might not
perhaps, for special equations, remain unchanged under still
other permutations.

If this case should occur, then the equation
(21) or (22) would have to have equal roots, or (20) equal or
opposite roots.
